%% GENERATED by tools/from_docs.py from stepss-docs. Do not edit by hand:
%% edit the page named below in stepss-docs and regenerate.
%% source: stepss-docs/src/content/docs/developer/uramses.mdx

\chapter{uRAMSES, the C and Fortran interface}\label{doc:uramses}

\href{https://github.com/SPS-L/stepss-uramses}{URAMSES} provides the framework and tools needed to compile and link custom Fortran models with stepss and the STEPSS GUI.

This is the route for user models. RAMSES resolves a model name against the
models compiled into it, so your own models are linked into the engine and
reach it the same way the built-in ones do. There is no way to load a model
library while a simulation is running.

\section{Prerequisites}\label{uramses:prerequisites}

\textbf{Linux (Debian, Ubuntu)}

\begin{itemize}
\tightlist
\item
  \textbf{gfortran} (GNU Fortran compiler)
\item
  \textbf{OpenBLAS} (optimized BLAS library)
\item
  \textbf{stepss}
\end{itemize}

\begin{Verbatim}
sudo apt install gfortran make libopenblas-dev
\end{Verbatim}

Debian and Ubuntu are the Linux distributions STEPSS is built and tested on; see
\hyperref[ch-install]{Installation}. The kit itself needs only
gfortran, GNU make and OpenBLAS, so it will build wherever those exist, but the
compiler has to match the ABI of the bundled \texttt{.mod} files and that is checked
against the distributions above.

\textbf{macOS}

\begin{itemize}
\tightlist
\item
  \textbf{gfortran} from Homebrew (macOS ships no Fortran compiler)
\item
  \textbf{OpenBLAS}, or Apple's Accelerate framework
\item
  \textbf{stepss}
\end{itemize}

\begin{Verbatim}
brew install gcc openblas
\end{Verbatim}

Apple Silicon only: the pre-compiled \texttt{modules\_m/} kit is arm64.

\textbf{Windows (MinGW)}

\begin{itemize}
\tightlist
\item
  \textbf{MSYS2} with the MinGW-w64 toolchain
\item
  \textbf{OpenBLAS}
\item
  \textbf{stepss}
\end{itemize}

\begin{Verbatim}
pacman -S mingw-w64-x86_64-gcc-fortran mingw-w64-x86_64-openblas
\end{Verbatim}

Build from the \textbf{MSYS2 MinGW 64-bit} shell. The plain MSYS shell links against
the \texttt{msys-2.0} runtime and produces a DLL that CPython cannot load.

\textbf{Windows (Intel)}

\begin{itemize}
\tightlist
\item
  \textbf{Microsoft Visual Studio} 2019 or later
\item
  \textbf{Intel oneAPI Fortran Compiler}, installed from Intel's own distribution;
  the \hyperref[ch-install]{Installation} page covers the gfortran
  toolchains only
\item
  \textbf{stepss}
\end{itemize}

\begin{docsnote}{Match gfortran to the kit}

Fortran \texttt{.mod} files can only be read by the gfortran release that wrote them.
Each kit records the exact module ABI version it was built with in its own
\texttt{BUILDINFO.txt} (\texttt{cat\ modules\_l/BUILDINFO.txt}), and the same information
appears in the toolchain table at the top of every URAMSES release. Run
\texttt{check-deps} on any Makefile to compare your compiler against the kit before
building; it fails early and names both versions on a mismatch.

If your distribution's default gfortran is the wrong generation, install a
matching one and pass it explicitly, e.g.
\texttt{make\ -f\ build/Makefile.linux\ FC=gfortran-11}.

\end{docsnote}

\section{Project Structure}\label{uramses:project-structure}

Every platform carries one suffix throughout: \texttt{l} for Linux, \texttt{m} for macOS,
\texttt{wg} for Windows/MinGW (gfortran) and \texttt{wi} for Windows/Intel. A module kit and
the output directory it builds into always share it.

\begin{Verbatim}
URAMSES/
+-- src/                         # Framework sources (all platforms)
|   +-- c_interface.f90          # C interface for Python integration
|   +-- main.f90                 # Main entry point (executable only)
|   +-- FUNCTIONS_IN_MODELS.f90  # Helper functions for models
|   +-- usr_exc_models.f90       # Exciter model associations
|   +-- usr_inj_models.f90       # Injector model associations
|   +-- usr_tor_models.f90       # Torque model associations
|   +-- usr_twop_models.f90      # Two-port model associations
|   +-- usr_dctl_models.f90      # Discrete controller associations
+-- custom_models/               # Your own models (all platforms)
|   +-- exc_*.f90                # Exciter models
|   +-- inj_*.f90                # Injector models
|   +-- tor_*.f90                # Torque models
|   +-- twop_*.f90               # Two-port models
|   +-- *.txt                    # Model parameter files
+-- modules_l/                   # Pre-compiled kit (Linux/gfortran)
|   +-- *.mod                    # Module interface files
|   +-- libramses.a              # Pre-compiled RAMSES library
+-- modules_m/                   # Pre-compiled kit (macOS arm64/gfortran)
|   +-- *.mod                    # Module interface files
|   +-- libramses.a              # Pre-compiled RAMSES library
+-- modules_wg/                  # Pre-compiled kit (Windows/MinGW gfortran)
|   +-- *.mod                    # Module interface files
|   +-- libramses.lib            # Pre-compiled RAMSES library
+-- modules_wi/                  # Pre-compiled kit (Windows/Intel Fortran)
|   +-- *.mod                    # Module interface files
|   +-- libramses.lib            # Pre-compiled RAMSES library
+-- build/                       # All build files
|   +-- Makefile.linux           # Linux builds
|   +-- Makefile.macos           # macOS builds (Apple Silicon)
|   +-- Makefile.windows         # Windows/MinGW builds
|   +-- msvs/                    # Visual Studio files (Windows/Intel)
|       +-- URAMSES.sln          # Visual Studio solution
+-- tools/                       # Kit-verification and release helpers
+-- Release_l/                   # Build output (Linux)
+-- Release_m/                   # Build output (macOS)
+-- Release_wg/                  # Build output (Windows/MinGW)
+-- Release_wi/                  # Build output (Windows/Intel)
\end{Verbatim}

Each platform links its own pre-compiled RAMSES kit, published with the matching
RAMSES release as \texttt{uramses-modules\_\{l,m,wg\}-\textless{}version\textgreater{}.zip}.

\begin{docsnote}{Note}

The Makefiles live in \texttt{build/} but are always run \textbf{from the repository root},
\texttt{make\ -f\ build/Makefile.linux}. \texttt{-f} does not change make's working directory,
so every path inside them stays root-relative. Do not \texttt{cd\ build}.

\end{docsnote}

\section{Model Types}\label{uramses:model-types}

{\def\LTcaptype{none} % do not increment counter
\begin{small}\begin{longtable}[]{@{}lll@{}}
\toprule\noalign{}
Type & Prefix & Description \\
\midrule\noalign{}
\endhead
\bottomrule\noalign{}
\endlastfoot
Exciters & \texttt{exc\_*} & Generator excitation system models \\
Injectors & \texttt{inj\_*} & Current/voltage injection models \\
Torque & \texttt{tor\_*} & Mechanical torque models \\
Two-port & \texttt{twop\_*} & Two-port network models (SVC, STATCOM, HVDC) \\
Discrete Control & \texttt{dctl\_*} & Discrete control system models \\
\end{longtable}\end{small}
}

\section{Building}\label{uramses:building}

\textbf{Linux}

Run from the repository root:

\begin{Verbatim}
# Build shared library + executable
make -f build/Makefile.linux all

# Build only the shared library (for stepss)
make -f build/Makefile.linux dll

# Build only the executable
make -f build/Makefile.linux exe

# Verify gfortran, OpenBLAS and the module kit
make -f build/Makefile.linux check-deps

# Clean build artifacts
make -f build/Makefile.linux clean
\end{Verbatim}

The shared library (\texttt{ramses.so}) will be in \texttt{Release\_l/}.

\textbf{macOS}

Run from the repository root:

\begin{Verbatim}
# Build shared library + executable
make -f build/Makefile.macos all

# Build only the shared library (for stepss)
make -f build/Makefile.macos dll

# Build only the executable
make -f build/Makefile.macos exe

# Verify gfortran, BLAS, the architecture and the module kit
make -f build/Makefile.macos check-deps

# Clean build artifacts
make -f build/Makefile.macos clean
\end{Verbatim}

The shared library (\texttt{ramses.so}) will be in \texttt{Release\_m/}.

Link against Apple's Accelerate framework instead of OpenBLAS with
\texttt{make\ -f\ build/Makefile.macos\ BLAS=accelerate}, and select a versioned Homebrew
compiler with \texttt{FC=gfortran-15}.

\begin{docsnote}{Note}

The macOS library is named \texttt{ramses.so}, not \texttt{ramses.dylib}. stepss separates
platforms using the \texttt{libs/lin}, \texttt{libs/mac} and \texttt{libs/win} folders rather than by
filename, so Linux and macOS share the same name.

\end{docsnote}

\textbf{Windows (MinGW)}

From the \textbf{MSYS2 MinGW 64-bit} shell, at the repository root:

\begin{Verbatim}
# Build shared library + executable
make -f build/Makefile.windows all

# Build only the shared library (for stepss)
make -f build/Makefile.windows dll

# Build only the executable
make -f build/Makefile.windows exe

# Verify gfortran, OpenBLAS and the module kit
make -f build/Makefile.windows check-deps

# Clean build artifacts
make -f build/Makefile.windows clean
\end{Verbatim}

The shared library (\texttt{ramses.dll}) will be in \texttt{Release\_wg/}.

This route is independent of the Intel Visual Studio projects; both can coexist
in the same clone, writing to \texttt{Release\_wg/} and \texttt{Release\_wi/}.

\begin{docsnote}{Note}

The Windows gfortran kit ships its archive as \texttt{libramses.lib} rather than
\texttt{libramses.a}. MinGW's linker does not probe that name for \texttt{-lramses}, so the
Makefile passes \texttt{modules\_wg/libramses.lib} by explicit path. Do not replace
that with \texttt{-lramses}.

\end{docsnote}

\begin{docsnote}{Note}

The resulting DLL depends on the MSYS2 runtime libraries (\texttt{libgfortran},
\texttt{libgcc\_s}, \texttt{libgomp}, \texttt{libopenblas}). Keep the MinGW \texttt{bin} directory on \texttt{PATH},
or copy those DLLs next to \texttt{ramses.dll}, when loading it from Python.

\end{docsnote}

\textbf{Windows (Intel)}

\begin{enumerate}
\def\labelenumi{\arabic{enumi}.}
\tightlist
\item
  Open \texttt{build/msvs/URAMSES.sln} in Visual Studio
\item
  Select the desired project:

  \begin{itemize}
  \tightlist
  \item
    \texttt{dllramses}: builds \texttt{ramses.dll} (for stepss)
  \item
    \texttt{exeramses}: builds \texttt{dynsim.exe} (standalone)
  \end{itemize}
\item
  Build the solution (Release x64 configuration)
\item
  Output will be in \texttt{Release\_wi/}
\end{enumerate}

\section{Adding Custom Models}\label{uramses:adding-custom-models}

\begin{enumerate}
\def\labelenumi{\arabic{enumi}.}
\item
  \textbf{Write the model source}: Create a \texttt{.f90} file in \texttt{custom\_models/} following the naming convention (\texttt{exc\_}, \texttt{inj\_}, \texttt{tor\_}, \texttt{twop\_}, or \texttt{dctl\_} prefix)
\item
  \textbf{Register the model}: Add the model association in the corresponding \texttt{usr\_*\_models.f90} file in \texttt{src/}
\item
  \textbf{Rebuild}: Compile and link the modified URAMSES with the Makefile for your
  platform (\texttt{build/Makefile.linux}, \texttt{build/Makefile.macos} or
  \texttt{build/Makefile.windows}), which picks up new \texttt{.f90} files in
  \texttt{custom\_models/} automatically. Only the Intel Visual Studio route needs the
  file added to the project by hand.
\end{enumerate}

\subsection{Example: Registering an Exciter Model}\label{uramses:example-registering-an-exciter-model}

In \texttt{src/usr\_exc\_models.f90}, declare the subroutine and point the model pointer at
it. The case label carries the \texttt{exc\_} prefix, because the router has already
added it to whatever name the data file used:

\begin{Verbatim}
   external :: exc_MY_EXCITER
   ...
   select case (modelname4)
      case('exc_MY_EXCITER')
         exc_ptr=>exc_MY_EXCITER
\end{Verbatim}

The other four routers follow the same shape, with \texttt{tor\_}, \texttt{inj\_}, \texttt{twop\_} and
\texttt{dctl\_} prefixes and their own pointer names. A data file then refers to the
model as either \texttt{MY\_EXCITER} or \texttt{exc\_MY\_EXCITER}.

Up to RAMSES 3.78 this was also how you reached a model that RAMSES compiled
but registered under no name, such as \texttt{exc\_AC8B} or \texttt{twop\_HVDC\_VSC}. From 3.79
every compiled model is registered, so that step is no longer needed. The
technique still works for anything you want to expose under a second name: the
subroutine is already in \texttt{libramses}, so no \texttt{external} declaration or source
file of your own is required.

\subsection{Included Example Models}\label{uramses:included-example-models}

The repository includes these example models:

{\def\LTcaptype{none} % do not increment counter
\begin{small}\begin{longtable}[]{@{}ll@{}}
\toprule\noalign{}
File & Description \\
\midrule\noalign{}
\endhead
\bottomrule\noalign{}
\endlastfoot
\texttt{exc\_ENTSOE\_lim.f90} & ENTSO-E exciter with limits \\
\texttt{inj\_AIR\_COND1\_mod.f90} & Air conditioning load injector \\
\end{longtable}\end{small}
}

\begin{docsnote}{Do not duplicate library models}

Models that the linked RAMSES library already provides, such as \texttt{exc\_ST1A},
\texttt{exc\_GENERIC} and \texttt{tor\_ENTSOE\_simp}, must not be copied into
\texttt{custom\_models/}. Defining a subroutine that the library also exports produces a
duplicate-symbol link error. Register those by name in the \texttt{usr\_*\_models.f90}
routers instead.

\end{docsnote}

\section{Using Custom Models}\label{uramses:using-custom-models}

\subsection{With stepss}\label{uramses:with-stepss}

Point stepss to your custom library:

\begin{Verbatim}
import stepss
ram = stepss.sim(custLibDir="path/to/Release_l")     # Linux
# ram = stepss.sim(custLibDir="path/to/Release_m")   # macOS
# ram = stepss.sim(custLibDir=r"path\to\Release_wg") # Windows (MinGW)
# ram = stepss.sim(custLibDir=r"path\to\Release_wi") # Windows (Intel)
\end{Verbatim}

stepss loads \texttt{ramses.so} on Linux and macOS, and \texttt{ramses.dll} on Windows.

\subsection{With STEPSS GUI}\label{uramses:with-stepss-gui}

The \textbf{Codegen} tab does this for you: it runs CODEGEN on your model
description, then drives these same Makefiles against the bundled URAMSES kit to
produce a custom \texttt{dynsim}, which it uses for subsequent simulations. You need
\texttt{gfortran}, GNU \texttt{make} and OpenBLAS installed; see
\hyperref[ch-install]{Installation}. Building by hand and replacing
the simulator manually is only necessary for models the Codegen tab cannot
generate.

\section{Helper Functions}\label{uramses:helper-functions}

The \texttt{FUNCTIONS\_IN\_MODELS.f90} module provides utility functions available in all
models. See the \hyperref[ch-libray_blocks]{Functions Reference} for the
complete list with signatures.

\section{Repository}\label{uramses:repository}

Source code: \href{https://github.com/SPS-L/stepss-uramses}{SPS-L/stepss-uramses}
